\section{The factorizable Whittaker category}
\label{section:Ran}

\subsubsection{} In the section we will review the definitions of the Ran version $\Whit(G)_\Ran$ of the Whittaker category from \cite{loc_glob_whit} and relate it to the configuration space version $\Whit(G)_\Conf$. Although $\Whit(G)_\Ran$ is a naturally a unital factorization category, we will not make explicit use of this additional structure.

\subsection{Review of Whittaker categories}

\subsubsection{Local model} For a finite set $I$, we write
    \[\Whit(G)_I = \DMod(\Gr_{G,I})^{\LL N,\psi}  \qquad  \Whit(G)_{I,\co} = \DMod(\Gr_{G, I})_{\LL N, \psi}\]
By \cite[Appendix B]{lin_nearby_cycles}, these assignments define factorization categories which we denote
  \[\Whit(G)  \qquad  \Whit(G)_\co.\]
Furthermore, we have a canonical $\DMod(X^I)$-linear equivalence
  \[\Whit_\psi(G)_I^\vee \simeq \Whit_{-\psi}(G)_{I,\co}\]
inducing an equivalence
  \[\Whit_\psi(G)^\vee \simeq \Whit_{-\psi}(G)_\co\]
of factorization categories.

\subsubsection{Local duality} The functor of renormalized $*$-averaging
    \[\Av_{*,\ren} : \DMod(\Gr_{G,I}) \to \Whit(\Gr_{G,I})\]
induces an equivalence
    \[\PsId_\Whit : \Whit(G)_{I,\co} \simto \Whit(G)_I.\]
We obtain an equivalence of factorization categories
    \[\Whit_{-\psi}(G)^\vee \simto \Whit_\psi(G).\]

\subsubsection{Global model} Write $\Bunc_{N,I}^\pol$ for the stack classifying a point $x \in X^I$, and a principal $G$-bundle on $X$ equipped with a degenerate $N^\omega$-reduction away from $x$. We set
    \[\Whit(G)_I^\glob = \Whit\left(\Bunc_{N,I}^\pol \right).\]
This defines a sheaf of categories over the Ran space, which we denote
    \[\Whit(G)^\glob.\]
Note that $\Whit(G)^\glob$ acquires a weak factorization structure from the weak factorization structure
  \[\Bunc_{N,\Ran}^\pol \underset{\Ran}{\times} \left[ \Ran \times \Ran \right]_\disj \to \left[ \Bunc_{N,\Ran}^\pol \times \Bunc_{N,\Ran}^\pol \right]_\disj.\]
The comparison with the local category will show that $\Whit(G)^\glob$ is a genuine factorization category.

\subsubsection{Global duality} The perfect pairing
    \[
    \Whit_{-\psi}\left( \Bunc_{N,I}^\pol \right) \otimes \Whit_{\psi}\left( \Bunc_{N,I}^\pol \right) \to \Vect
    \qquad \langle \WW_1, \WW_2 \rangle = \Gamma\left( \Bunc_{N,I}^\pol, \WW_1 \shimes \WW_2 \right)\]
defines a factorizable equivalence
  \[\DD : \Whit_{\psi}(G)^\glob \simto \Whit_{-\psi}(G)^{\glob,\vee}.\]

\subsubsection{Local v global} The map
    \[\pi : \Gr_{G,I} \to \Bunc_{N,I}^{\pol}\]
induces equivalences
    \[
        \pi^! : \Whit(G)_I^\glob \to \Whit(G)_I  \qquad
        \pi_* : \Whit(G)_{\co,I} \to \Whit(G)_{\co, I}^\glob.
    \]
We use the same notation for the Ran versions. The interaction of this comparison with duality is summarized by the diagram below.
\begin{equation}
\label{eqn:Whit_duality_loc_global}
\begin{tikzcd}[
      column sep={5em,between origins},
      row sep={6em,between origins}
    ]
    \Whit_{-\psi}(G)^{\glob,\vee} & & & \Whit_\psi(G)^\glob \arrow[lll, "\DD" above] \arrow[dr, "\pi^!"] & \\
    \Whit_{-\psi}(G)^\vee \arrow[rr, "\sim"] \arrow[u, "(\pi^!)^\vee"] & & \Whit_\psi(G)_\co \arrow[rr, "\PsId{[-2 \delta_N]}"] \arrow[ur, "\pi_*"] & & \Whit_\psi(G)
\end{tikzcd}
\end{equation}

\subsubsection{Hecke action} 
\begin{itemize}
    \item We write $\Bun_{G,I}$ for $\Bun_G \times X^I$ and $\DMod(\Bun_G)_I$ for $\DMod(\Bun_{G,I})$, viewed as a $\DMod(X^I)$-module.
    \item Recall the Hecke stack
        \[
        \begin{tikzcd}[
      column sep={5em,between origins},
      row sep={2em,between origins}
    ]
          & \HH_{G,I} \arrow[dl, "\hl" above left] \arrow[dr, "\hr"] \arrow[dd, "s"] &    \\
          \Bun_G & & \Bun_G \\
          & X^I &
        \end{tikzcd}
        \]
    Note that $\DMod(\HH_{G,I})$ is naturally a $\DMod(X^I)$-algebra acting on $\DMod(\Bun_G)_I$.
    \item We write $\Sph_{G,I}$ for the spherical Hecke category over $X^I$. It is a $\DMod(X^I)$-algebra acting on $\DMod(\Bun_G)_I$ through the canonical monoidal functor
        \[
        \Sph_{G,I} \to \DMod(\HH_{G,I})   
        \qquad \SS \mapsto \HH(\SS).
        \]
    \item We write $\Sph_{G,\Ran}$ for the Ran version of $\Sph_{G,I}$. It is a monoidal factorization category.
\end{itemize}
Now $\Sph_{G,I}$ acts on all of the previously defined Whittaker categories over $X^I$. Let us discuss the compatibility of (\ref{eqn:Whit_duality_loc_global}) with the $\Sph_G$ action. The right triangle of (\ref{eqn:Whit_duality_loc_global}) is a diagram of right $\Sph_G$-modules, and the left arrow $(\pi^!)^\vee$ is a morphism of left $\Sph_G$-modules. The global Verdier duality is $\Sph_G$-linear with respect to the canonical anti-involution $\tau : \Sph_G \to \Sph_G^\rev$, as is the identification $\Whit_{-\psi}(G)^\vee \simto \Whit_\psi(G)_\co$.

\subsubsection{Whittaker coefficients and Poincare series} The functor of Whittaker coefficients
    \[
    \coeff_{I,*} : \DMod(\Bun_G)_I \overset{\pi^!}{\longrightarrow} \DMod\left(\Bunc_{N,I}^\pol \right) \overset{\Av_*}{\longrightarrow} \Whit(G)_I^\glob
    \]
is $\Sph_{G,I}$-linear. Therefore, its left adjoint
    \[\Poinc_{I,!} : \Whit(G)_I^\glob \to \DMod(\Bun_{G,I})\]
is automatically co-lax $\Sph_{G,I}$-linear. We write $\coeff_{\Ran,*}$ and $\Poinc_{\Ran,!}$ for the Ran variants of these functors.

\begin{proposition}
\label{proposition:Poinc_linear}
The functor $\Poinc_{I,!}$ is actually $\Sph_{G,I}$-linear.
\end{proposition}

\subsubsection{A triviality} We will use the following observation in both the proof of Proposition \ref{proposition:Poinc_linear} and in later sections without comment. Suppose that $f : M_1 \to M_2$ is a colax morphism of modules over a monoidal category $A$. Let us say that an object $m$ of $M_1$ is \emph{strict} if the colax constraint
  \[f(a \otimes m) \to a \otimes f(m)\]
is an isomorphism for all $a \in A$. Then the full subcategory of strict objects is a monoidal ideal.

\subsubsection{Proof of Proposition \ref{proposition:Poinc_linear}} Since $\Whit(G)_I$ is generated as a $\Sph_{G,I}$-module by the vacuum object
    \[\WW_{\vac, I} = \WW_\vac \boxtimes \omega_{X^I} \in \Whit\left(\Bunc_N \times X^I\right) \mon \Whit(G)_I,\]
it suffices to show that the colax constraint
    \[\Poinc_{I,!} \cdot (\SS \star \WW_{\vac,I})
    \to \SS \star \left( \Poinc_{I,!} \cdot \WW_{\vac,I} \right)\]
is an isomorphism. Now consider the following diagram.
\[
\begin{tikzcd}[
      column sep={4em,between origins},
      row sep={4em,between origins}
    ]
    & \Bunc_N \!\underset{\Bun_G}{\times}\! \HH_{G,I} \arrow[dl, "\hl" above left] \arrow[dr] \arrow[drrr, bend left=15, dashed, "\hr_{\!N}"] & & & \\
  \Bunc_N \arrow[dr, "\pi"] & & \HH_{G,I} \arrow[dl, "\hl"] \arrow[dr, "\hr \times s"] & & \Bunc_{N,I}^\pol \arrow[dl, "\pi_I"] \\
    & \Bun_G & & \Bun_{G,I} & 
\end{tikzcd}
\]
The colax constraint identifies with
\[
    \pi_{I,!} \cdot \hr_{\!N,!} \cdot \left( \WW_\vac \shoximes \SS_\HH \right)
    \to (\hr \times s)_* \cdot \left( \hl^! \cdot \pi_! \cdot \WW_\vac \shimes \SS_\HH \right),
\]
which is an isomorphism because $\hr \times s$ and $\hr_{\!N}$ are proper and $\SS_\HH$ is ULA over $\Bun_G$.

\subsubsection{} We write
    \[\Gamma(\Ran, -) : \DMod(\Bun_G)_\Ran \to \DMod(\Bun_G)\]
for pushforward along the projection
    \[\pr : \Bun_G \times \Ran \to \Bun_G.\]
It is left adjoint to $\pr^!$, and we have a commutative diagram
\[
\begin{tikzcd}[
      column sep={9em,between origins},
      row sep={4em,between origins}
    ]
    \DMod(\Bun_G)_\Ran^\vee \arrow[r, "\PsId_!"] \arrow[d, "(\pr^!)^\vee" left] & \DMod(\Bun_G)_\Ran \arrow[d, "{\Gamma(\Ran,-)}"] \\
    \DMod(\Bun_G)^\vee \arrow[r, "\PsId_!"] & \DMod(\Bun_G)
\end{tikzcd}
\]

\subsection{Conf v Ran}

\subsubsection{} The purpose of this subsection is to construct a functor
    \[\sprd_* : \Whit(M)_\Conf \to \Whit(M)_\Ran\]
along with the following compatibilities:
\begin{itemize}
    \item We have
        \begin{equation}
        \label{eqn:Poinc_Ran_sprd}
        \Gamma(\Ran, \Poinc_{\Ran,!} \cdot \sprd_*) \simeq \Poinc_{\Conf, !}.
        \end{equation}
    \item Let $\sprd^!$ denote the right adjoint to $\sprd_*$. Then there is a commutative diagram:
    \begin{equation}
    \label{eqn:sprd_duality}
    \begin{tikzcd}[
      column sep={9em,between origins},
      row sep={4em,between origins}
    ]
      \Whit(M)_\Conf^\vee \arrow[r, "\sim"] \arrow[d, "(\sprd^!)^\vee"] & \Whit(M)_\Conf \arrow[d, "\sprd_*"] \\
      \Whit(M)_\Ran^\vee \arrow[r, "\sim"] & \Whit(M)_\Ran
    \end{tikzcd}
    \end{equation}
\end{itemize}

\subsubsection{} The device relating $\Conf$ to $\Ran$ is a Cartesian diagram
\begin{equation}
\label{eqn:Conf_Ran_vac}
\begin{tikzcd}[
  column sep={7.5em,between origins},
  row sep={2em,between origins}
]
    & \Bunc_{N_M, \Ran}^{\pol, \neg} \arrow[dl, "\aph\pi" above left] \arrow[dr, "\aph i^\neg"] \arrow[dd] & \\
  \Bunc_{N_M, \Conf}^{\pol} \arrow[dd] & & \Bunc_{N_M, \Ran}^{\pol} \arrow[dd] \\
    & \Gr_{M^\ab, \Ran}^{\omega,\neg} \arrow[dl, "\pi" above left] \arrow[dr, "i^\neg"] &  \\
  \Conf_{G,P} & & \Gr^\omega_{M^\ab, \Ran}
\end{tikzcd}
\end{equation}
whose terms are defined as follows:
\begin{itemize}
    \item We write $\Gr_{M^\ab,\Ran}^\omega$ for the $\omega^\rhov$-twisted form of the affine Grassmannian. It classifies a point $x$ of $\Ran$, a principal $M^\ab$-bundle $\FF_{M^\ab}$, and a system of compatible maps
        \[
        \FFo_{M^\ab} \overset{M^\ab}{\times} \theta \to \omega^{\theta(\rhov)}
        \qquad
        \theta \in \Lambda_{G,P}.
        \]
    We write $\Gr_{M^\ab}^{\omega,\neg}$ for the closed sub-prestack where these maps extend to regular maps for all $\theta \in \Lambda_{G,P}^\pos$. When this happens we obtain a divisor $D \in \Conf_{G,P}$ supported on $x$. This defines the leftward map denoted $\pi$.
    \item The prestack $\Bunc_{N_M, \Ran}^\pol$ is the Ran version of $\Bunc_{N_M,I}^\pol$. It classifies a point $x$ of the Ran space, a principal $M$-bundle $\FF_M$ on $X$, and an $N_M^\omega$-Pl\"ucker data away from $x$. The right vertical map is defined by observing that such a Pl\"ucker data provides an identification between the bundles
        \[
        \FF_M \overset{M}{\times} M^\ab
        \qquad  \text{and} \qquad
        \omega^\rhov \overset{T}{\times} M^\ab
        \]
    away from $x$.
    \item We may define $\Bunc_{N_M, \Ran}^{\pol, \neg}$ as the fiber product of either face; the two are canonically identified.
\end{itemize}
Then we define
  \[\sprd_* = (\aph i^\neg)_* \cdot \aph\pi^!.\]

\subsubsection{} 
\label{subsection:sprd_compatibilities}
In fact, the map
    \[\Gr_{M^\ab}^{\omega,\neg} \mon \Conf \times \Ran\]
identifies $\Gr_{M^\ab}^{\omega,\neg}$ with the locus of pairs $(D, x)$ cut out by the closed condition $|D| \mon x$. Let us use this to obtain the desired compatibilities.
\begin{itemize}
    \item Since $\pi$ is pseudo-ind-proper and $i^\neg$ is a closed immersion, $\sprd_*$ admits a continuous right adjoint
        \[\sprd^! \simeq \aph\pi_* \cdot (\aph i^\neg)^!.\]
    This gives the compatibility (\ref{eqn:sprd_duality}) with duality.
    \item The compatibility (\ref{eqn:Poinc_Ran_sprd}) with Poincare series follows from the universal homological contractibility of $\pi$.
\end{itemize}